\chapter{Conclusion}
\label{ch:conclusion}

\section{Summary of contributions}

This thesis developed two approaches for learning macroscopic structure
from noisy interacting-agent trajectories: equation-free
forecasting~\cite{alvarez2025eqfree} via shift-aligned POD with latent
MVAR and LSTM time-steppers (Chapter~\ref{ch:rom_results}), and
weak-form mechanistic identification via
WSINDy~\cite{messenger2021wsindy} (Chapter~\ref{ch:wsindy_results}).
The two branches yield complementary diagnostics: regimes where
identification succeeds but autonomous rollout fails expose closure
limitations that neither branch could reveal alone.

\emph{Part~I} established an equation-free forecasting baseline.
Chapter~\ref{ch:background} laid out the Vicsek--Morse microscopic model
and KDE coarse-graining.
Chapters~\ref{ch:global_pod_pipeline}--\ref{ch:latent_rom} developed the
shift-aligned POD compression and the MVAR/LSTM latent time-steppers.
Chapter~\ref{ch:exp_protocol_metrics} defined the nine-regime benchmark
and evaluation metrics, and Chapter~\ref{ch:rom_results} delivered the
main forecasting result: MVAR achieves $R^2>0.99$ on the strongest
constant-speed regimes while requiring $116\times$ fewer parameters than
LSTM, confirming that aligned latent dynamics are effectively linear for
well-shifting regimes.  Variable-speed regimes exposed a
shift-predictability bottleneck that neither model class could overcome.

\emph{Part~II} turned from forecasting to mechanism.
Chapter~\ref{ch:wsindy_methodology} formulated the coupled multi-field
WSINDy problem, with a physically motivated Morse library, automatic
support-scale selection, and MSTLS sparsification.
Chapter~\ref{ch:wsindy_results} applied this pipeline to all nine
regimes, achieving density weak-form $R^2>0.99$ in eight of nine cases
and recovering the continuity backbone universally.  The CS/VS paired
comparison showed that variable speed degrades not only rollout but
also momentum identification: the three-field library
$(\rho,p_x,p_y)$ is structurally insufficient to close the mean-field
description when agent speed is state-dependent, constituting a
structural non-identification rather than a regression failure.

Across the benchmark regimes, the main empirical pattern was clear:
density weak-form $R^2$ exceeded~0.99 on every regime except
blackhole\_VS (Table~\ref{tab:mechanism_confidence}), with
the continuity backbone recovered in all nine cases;
the EF-ROM achieved forecast $R^2>0.99$ on the strongest constant-speed
regimes (CS\,+\,PV median $0.994$), while VS median forecast $R^2$
dropped to $0.425$ for MVAR and $-0.052$ for LSTM; and WSINDy
nevertheless recovered sparse regime-consistent effective mechanisms even
where autonomous rollout was poor.

Weak-form sparse regression does detect interpretable coarse-grained
structure in these data.  The remaining challenge---the \emph{closure gap}---is whether
the selected state variables and library are sufficiently rich to close
that structure under autonomous evolution at the chosen effective scale.
Bridging this gap is a structural problem, not primarily a fitting
problem: it requires richer effective states or constrained regression
that enforces stability by construction.

\section{Limitations}
\label{sec:limitations}

These conclusions hold for a specific effective-modeling pipeline and should
be interpreted subject to four main limitations.

\paragraph{Scale and observability.}
The discovered PDE is scale-dependent.  KDE bandwidth, weak test-function
support, density transformation, and translational alignment together define
what structures are observable and therefore identifiable.  The recovered law
is thus an effective finite-scale description, not a unique continuum truth
independent of coarse-graining.  This matters especially for short-range
interaction effects, rotational organisation, and multimodal transport,
which are only partially resolved by the present preprocessing choices.

\paragraph{Finite-$N$ interpretation.}
With $N{=}300$ agents, the density
field is a noisy estimator of the mean-field limit, so the WSINDy
coefficients should be interpreted as effective finite-$N$ parameters
rather than as rigorous continuum-limit PDE coefficients.  These particle
counts reflect a practical compute constraint: each regime requires
approximately 2--4~CPU-hours on the Oscar HPC cluster; scaling to
$N{=}1000$ increases the per-regime cost roughly tenfold, which exceeded
the available allocation for a full nine-regime benchmark.
A preliminary $N$-convergence study was conducted, but a full cross-regime
convergence analysis remains open.

\paragraph{Closure and task asymmetry.}
The main limitation of the WSINDy branch is that the selected state
variables and candidate library may not close the effective law at the
chosen scale (Section~\ref{sec:closure_problem}).  The EF-ROM branch
bypasses this by learning a closed discrete-time map in latent space,
so side-by-side metrics must be interpreted accordingly.

\paragraph{Coefficient uncertainty.}
No bootstrap confidence intervals are reported for the WSINDy
coefficients.  Cross-scale stability selection
(Section~\ref{subsec:stability}) serves as a proxy for
robustness, but a formal uncertainty quantification campaign---e.g.\
resampling over initial-condition draws or KDE bandwidth
perturbations---would strengthen the coefficient-level conclusions.

\paragraph{Generalisability.}
Even with nine benchmark regimes and four initial-condition families, the
study remains restricted to one simulator family on a periodic flat domain.
The conclusions should therefore be read as claims about the explored
Vicsek--Morse~\cite{vicsek1995vicsek,dorsogna2006morse} regime family under this coarse-graining pipeline, not as
universal statements about collective-motion PDE discovery.

\paragraph{Threats to validity.}
Five additional experimental factors may restrict generalisability or bias the
reported conclusions.
(i)~All data come from one Vicsek--Morse model on a periodic square
domain at $N{=}300$; identified structures may not transfer to
larger swarms, other geometries, or alternative interaction models.
(ii)~Both branches consume density fields from the same fixed-bandwidth
Gaussian KDE, so any coarse-graining bias propagates to every subsequent
result, making the two branches less independent than they appear.
(iii)~The WSINDy candidate library is designed with knowledge of the
expected continuum mechanisms (advection, diffusion, alignment), so the
regression cannot discover mechanisms absent from the library.
(iv)~POD rank, MVAR lag, test-function support, and thresholding schedule
are held fixed across all nine regimes, which may disadvantage regimes
whose dynamics would benefit from individual tuning.
(v)~Shift alignment assumes a single dominant translational mode; regimes
with rotational or multi-modal transport violate this assumption, inflating
both ROM error and WSINDy closure residuals.

\section{Future work}
\label{sec:future_work}

Six extensions follow most directly from the thesis's findings.

\emph{1.~Richer closure and state augmentation.}
Closure insufficiency is the main barrier to autonomous PDE rollout
(Section~\ref{sec:closure_problem}).  The most direct next step is therefore
to enlarge the effective state or the candidate library, for example through
congestion/saturation terms, higher-order moment surrogates, or additional
fields that capture variability unresolved by $(\rho,\mathbf{p})$ alone.

\emph{2.~Constrained sparse regression.}
The closure gap could be narrowed from the regression side by replacing
unconstrained MSTLS with a stability-aware formulation such as SR3 (Sparse
Relaxed Regularised Regression)~\cite{zheng2019sr3} or by imposing physical
constraints during optimisation (for example, forcing the Laplacian
coefficient to be non-negative so that the identified diffusion operator
is dissipative.  Such {\em physical inductive biases} would convert
post-hoc safeguards (Section~\ref{sec:pipeline_fixes}) into structural
guarantees, reducing the gap between identification quality and rollout
stability.

\emph{3.~Scale-sensitivity analysis.}
The KDE bandwidth~$\sigma$ acts as a spatial filter: every identified PDE
is implicitly a function of that coarse-graining scale.  A systematic
Pareto-front analysis trading KDE bandwidth~$\sigma$ (smaller~$\to$
finer spatial resolution) against weak test-function support
(larger~$\to$ more noise averaging) would
map how the identified law changes with coarse-graining and would quantify
which interaction mechanisms are robustly recoverable at which scales.
This extension would clarify whether the recovered coefficients are
properties of the physics or of the smoothing---the most immediate
methodological question raised by the present results.

\emph{4.~Adaptive POD rank.}
The fixed rank $m{=}19$ used throughout this thesis may undersample the
latent dynamics of variable-speed regimes, where spectral decay is slower.
Future work should replace the fixed rank with an energy-based criterion
(e.g.\ retaining ${\geq}\,99.9\%$ of cumulative energy) or with the
singular value hard threshold of Gavish and
Donoho~\cite{gavish2014svht}, which provides an asymptotically optimal
cutoff for matrices with known noise level, allowing the ROM dimension
to adapt to regime complexity.

\emph{5.~Finite-$N$ and out-of-family validation.}
A broader $N$-convergence study, together with tests on additional simulator
families or geometries, is needed to determine which identified structures
persist beyond the present finite-$N$, single-simulator setting.

\emph{6.~Extension to curved domains.}
All numerical results in this thesis use flat periodic domains, but the
weak-form integral framework extends naturally to Riemannian manifolds
$(\mathcal{M},g)$: test functions are defined on the manifold's intrinsic
coordinates, and the integration-by-parts identity that eliminates spatial
derivatives replaces the Euclidean divergence theorem with its Riemannian
counterpart~\cite{lee2018riemannian}.  Concretely, the weak-form residual
for a conservation law $\partial_t u + \nabla_g\!\cdot\!\mathbf{F}=0$
becomes
\begin{equation}
  \int_{\mathcal{M}} \partial_t u\;\varphi\;\mathrm{d}V_g
  \;=\;
  \int_{\mathcal{M}} \mathbf{F}\cdot\nabla_g\varphi\;\mathrm{d}V_g,
  \label{eq:weak_form_riemannian}
\end{equation}
where $\nabla_g$ is the Levi-Civita connection, $\mathrm{d}V_g =
\sqrt{\det g}\;\mathrm{d}x^1\!\wedge\!\cdots\!\wedge\mathrm{d}x^n$ is
the Riemannian volume form, and the boundary term vanishes for compactly
supported~$\varphi$.  Two additional requirements arise on curved
domains.  First, all library terms must be \emph{metric-aware}: spatial
derivatives become covariant derivatives, and Laplacians are replaced by
the Laplace--Beltrami operator $\Delta_g u = \mathrm{div}_g(\nabla_g u)$.
Second, any time-stepping scheme that updates fields in a tangent-space
linearisation must \emph{project back to the manifold} at each step,
for example via the Riemannian exponential map
$u_{k+1} = \exp_{u_k}(\Delta t\;\dot{u}_k)$, to ensure that the
evolved density remains a well-defined field on~$\mathcal{M}$.
The main implementation task is
replacing FFT-based convolutions (used here for the Morse potential
$\Phi=W*\rho$) with manifold-aware kernel evaluations, for example via
heat-kernel approximations or spherical-harmonic transforms on~$S^2$.
This extension would enable WSINDy-based PDE discovery for agent-based
models on curved surfaces---a setting increasingly relevant in
biological collective motion (cells on tissue surfaces) and
geophysical swarming (flocks navigating the Earth's curvature).
The repository accompanying this thesis is named
\texttt{wsindy-manifold} in anticipation of this direction.

\noindent
Taken together, the equation-free and equation-based branches define a
\emph{trade-off map} for data-driven modelling of agent-based systems.
The EF-ROM delivers high short-term forecasting accuracy with zero
physical interpretability; WSINDy delivers mechanistic insight with
limited autonomous forecasting capability (due to the closure gap).
This work demonstrates that we can successfully ``read'' the physics of
interacting-agent systems from data; ``writing'' a stable predictive
simulation from those discovered laws remains the primary challenge for
the next generation of data-driven modelling.
